% Ciclo di vita di una transazione Bitcoin, disegnato in TikZ.
% Caricato nel preambolo; definisce il comando
%   \ciclovita{<fase da evidenziare: 1-4, 0 = nessuna>}{<testo della nota, vuoto = nessuna nota>}
\definecolor{cvevidenza}{HTML}{C00000}
\definecolor{cvnota}{HTML}{E7E6E6}
\definecolor{cvnotabordo}{HTML}{7F7F7F}

\newcommand{\ciclovita}[2]{%
\scalebox{0.85}{% dimensione della figura
\begin{tikzpicture}[
    icona/.style={font=\fontsize{30}{30}\selectfont},
    intestazione/.style={font=\Large},
    numero/.style={fill=black, text=white, font=\Large, minimum width=0.75cm,
        minimum height=0.95cm, inner sep=0pt},
    fase/.style={align=left, font=\large, anchor=west, inner sep=2pt},
    freccia/.style={-{Triangle[length=5mm, width=6mm]}, line width=2.5pt}
]
    % attori
    \node[icona]        at (1.2, 5.2)  {\faUser};
    \node[intestazione] at (2.9, 5.2)  {Alice};
    \node[icona]        at (1.2, 3.9)  {\faCoins};
    \node[icona]        at (2.9, 3.9)  {\faSignature};
    \node[intestazione] at (8.3, 5.2)  {Network};
    % rete peer-to-peer: nodi collegati fra loro
    \begin{scope}[shift={(8.3, 3.95)}]
        \coordinate (r1) at (-0.65, 0.05);
        \coordinate (r2) at (-0.1, 0.5);
        \coordinate (r3) at (0.6, 0.35);
        \coordinate (r4) at (0.35, -0.4);
        \coordinate (r5) at (-0.35, -0.5);
        \draw[line width=1.6pt] (r1) -- (r2) -- (r3) -- (r4) -- (r5) -- (r1)
            (r2) -- (r4) (r1) -- (r4);
        \foreach \r in {r1, r2, r3, r4, r5}
            \fill (\r) circle (0.16);
    \end{scope}
    \node[intestazione] at (14.4, 5.2) {Blockchain};
    \node[icona]        at (14.4, 3.9) {\faCubes};
    \draw[freccia] (4.4, 4.2) -- (6.4, 4.2);
    \draw[freccia] (10.2, 4.2) -- (12.2, 4.2);

    % fasi
    \node[numero] (n1) at (0.4, 2.4)  {1};
    \node[fase]   (f1) at (0.9, 2.4)  {Creazione e firma\\della transazione};
    \node[numero] (n2) at (6.2, 2.4)  {2};
    \node[fase]   (f2) at (6.7, 2.4)  {Broadcast e validazione\\della transazione};
    \node[numero] (n3) at (12.3, 2.4) {3};
    \node[fase]   (f3) at (12.8, 2.4) {Mining o creazione\\dei blocchi};
    \draw[freccia] (14.4, 1.6) -- (14.4, 0.35);
    \node[icona]       at (14.4, -0.4) {\faCubes};
    \node[numero] (n4) at (10.2, -0.4) {4};
    \node[fase]   (f4) at (10.7, -0.4) {Transazione\\confermata};

    % fase evidenziata
    \ifnum#1>0
        \node[draw=cvevidenza, line width=1.3pt, inner sep=5pt, fit=(n#1)(f#1)] {};
    \fi

    % nota
    \if\relax\detokenize{#2}\relax\else
        \node[draw=cvnotabordo, fill=cvnota, anchor=west, text width=7.6cm,
            align=left, font=\large, inner sep=6pt] at (0, -0.4) {#2};
    \fi
\end{tikzpicture}}%
}
